\chapter{Conclusion}
\label{sec:conclusion}

This thesis asked under what conditions firm-level AI investments translate into
competitive advantage and superior firm performance, and answered it with a systematic
review of 67 empirical studies published between 2020 and 2026, drawn from a documented
Scopus search of 432 records, screened in three recorded stages, and coded by two
independent coders against a frozen scheme. The coding table in the appendix is the
review's core deliverable; every number in the text derives from it.

The answer favors conditions over the investment itself. Across the corpus, 38 of the
67 studies report positive average effects, and the evidence is conditional almost
throughout: 63 studies identify at least one condition
that decides whether the investment reaches its outcome. The review ordered these
conditions into six groups, from the complementary investments a firm makes alongside
the technology to the credibility with which the investment is communicated. For
competitive advantage, the evidence is thin. 15 studies measure the construct, almost
always as a perceived edge at one point in time, and only a single study in the corpus
follows an advantage against imitation over time.

The contribution is the ordered map itself. Where individual studies report single
moderators and mediators, the review shows that the conditions recur in patterned
groups across methods and settings, that the configurational evidence finds them
working in combination rather than alone, and that the two outcome constructs the field
often blends carry different amounts of evidence. The map marks where the next studies
are needed, with persistence the least covered and most consequential claim. Firms can
read the same map as a checklist of what has to be in place before the investment can
be expected to pay.
